% English body fragment for SGA 2 XII (en-4.tex)

\section[Universal finiteness conditions]
{Universal finiteness conditions for a non-proper morphism} \label{XII.5} Let us recall for the record the \setcounter{equation}{25}

\begin{proposition} \label{XII.5.1}
Let $f:X\to S$ be a proper morphism of preschemes, with $S$ locally noetherian, $U$ an open subset of $X$, $g:U\to X$ the canonical immersion, $h=fg:U\to S$, $F$ a Module on $U$. Suppose that the Modules $\R^ig_\ast(F)$ are coherent for $i\leq n$ (hypothesis of a local nature on $X$, which is verified in practice with the help of the criterion \Ref{VIII}~\Ref{VII.2.3}). Then $\R^ih_\ast(F)$ is coherent for $i\leq n$.
\end{proposition}

This follows at once from the Leray spectral sequence
\steco
\begin{equation} \label{eq:XII.26} {} \E_2^{p, q}=\R^pf_\ast(\R^qg_\ast(F)) \To \R^\ast h_\ast(F),
\end{equation}
and from the fact that the higher direct images under $f$ of a coherent Module on $X$ are coherent (\EGA III 3.2.1).

\begin{proposition} \label{XII.5.2}
Let $S$ be a locally noetherian prescheme, $\Ss$ a quasi-coherent graded Algebra, of finite type over $S$, generated by $\Ss_1$, $X$ a subprescheme of $\Proj(\Ss)$, $\OX(1)$ the invertible Module on $X$ very ample relative to $S$ induced by $\Proj(\Ss(1))$, $U$ an open subset of $X$, $g:U\to X$ the canonical immersion, $h=fg:U\to S$, $F$ a quasi-coherent Module on $U$, whence twisted Modules $F(m)=F\otimes \OX(m)$) ($m\in\ZZ$), $n$ an integer, $m_0$ an integer. The following conditions are equivalent:
\begin{enumeratei}
\item
$\R^ig_\ast(F)$ is coherent for $i\leq n$.
\item
$\bigoplus_{m>m_0}\R^ih_\ast(F(m))$ is an $\Ss$-Module of finite type for $i\leq n$.
\end{enumeratei}
\end{proposition}

\begin{proof}
Replacing in the spectral sequence above $F$ by $F(m)$ one
finds a spectral sequence of graded $\Ss$-Modules
\[
\E_2^{p, q}= \tbigoplus_{m\geq m_0} \R^pf_\ast(\R^qg_\ast(F(m)) \To \tbigoplus_{m\geq m_0} \R^\ast h_\ast(F(m)).
\]
As one has
\[
\R^qg_\ast(F(m)) \simeq \R^qg_\ast(F) (m),
\]
one sees that if the $\R^ig_\ast(F)$ are coherent, $\E_2^{p, q}$ is of finite type over $\Ss$ for $q\leq n$, thanks to part a) of lemma \Ref{XII.5.3} below, which implies that the abutment is of finite type over $\Ss$ in degree $i\leq n$. This proves (i) \ALORS (ii). Moreover, reasoning in the abelian category of graded $\Ss$-Modules modulo the thick subcategory $\mathcal{C}$ of those which are quasi-coherent of finite type, one finds by the preceding spectral sequence
\[
\tbigoplus_{m\geq m_0} \R^{n+1} h_\ast(F(m)) \simeq \tbigoplus_{m\geq m_0} f_\ast (\R^{n+1} g_\ast(F)(m)) \mod \mathcal{C},
\]
which proves that if the left-hand side is an $\Ss$-Module of finite type, then $\R^{n+1} g_\ast(F)$ is coherent, by virtue of part b) of lemma \Ref{XII.5.3}. This proves the implication (ii) \ALORS (i) by induction on $n$. It remains to prove:

\begin{lemma} \label{XII.5.3} Let $S$, $\Ss$, $X$, $f$ be as in \Ref{XII.5.2}, and $G$ a quasi-coherent Module on $X$, $m_0$ an integer. Then
\begin{enumeratea}
\item
If $G$ is coherent, then for every integer $i$, the graded Module
\[
\tbigoplus_{m\geq m_0} \R^i f_\ast(G(m))
\]
over $\Ss$ is of finite type.
\item
Conversely, suppose that the Module $\bigoplus_{m\geq m_0} \R^i f_\ast(G(m))$ over $\Ss$ is of finite type, then $G$ is coherent.
\end{enumeratea}
\end{lemma}

\begin{proof}[Proof of \Ref{XII.5.3}]
For a), the case $i=0$ is given in \EGA III 2.3.2, the case $i>0$ in \EGA III 2.2.1 (i)(ii) which says that the $ \R^i f_\ast G(m)$ are coherent, and vanish for $m$ large (if one assumes $S$ noetherian, which is permissible). For b), one notes that $G$ is isomorphic to $\SheafProj (\bigoplus_{m\geq m_0}f_\ast(G(m))$ (\EGA II 3.4.4 and 3.4.2), which proves that $G$ is coherent if $\bigoplus_{m\geq m_0}f_\ast(G(m))$ is of finite type over $\Ss$, by virtue of \loccit 3.4.4.
\skipqed
\end{proof}
\skipqed
\end{proof}

\begin{corollary}[{{\normalfont$S$ noetherian}}] \label{XII.5.4}
Suppose that $\R^ig_\ast(F)$ is coherent for $i\leq n$, then for $i\leq n+1$, and $m$ large, one~has a canonical isomorphism:
\[
\R^ih_\ast(F(m)) \simeq f_\ast(\R^ig_\ast(G)(m)).
\]
\end{corollary}

Indeed the spectral sequence (\Ref{eq:XII.26}) for $F(m)$ then degenerates in degree $\leq n$, by \EGA III 2.2.1~(ii), whence at once the result (which moreover recovers the implication (ii) \ALORS (i) of~5.2).

%
\begin{corollary} \label{XII.5.5}
Under the preliminary conditions of \Ref{XII.5.2}, $S$ noetherian, the following conditions are equivalent:
\begin{enumeratei}
\item
$\bigoplus_{m\geq m_0}h_\ast(F(m)$ is of finite type over $S$, and $\R^ih_\ast(F(m))=0$ for $0<i\leq n$ and $m$ large.
\item
$g_\ast(F)$ is coherent, and $\R^ig_\ast(F)=0$ for $0<i\leq n$.
\item[\textup{(ii bis)}] $g_\ast(F)$ is coherent, and $\prof_Y g_\ast(F)>n+1$.
\end{enumeratei}
\end{corollary}

The equivalence of (ii) and (ii bis) is contained in \Ref{III}~\Ref{III.3.3}. Moreover, by virtue of \Ref{XII.5.2} the conditions (i) and (ii) both imply that the $\R^ig_\ast(F)$ ($i\leq n$) are coherent. The equivalence of (i) and (ii) then follows from \Ref{XII.5.4}, taking into account the fact that every a coherent Module $G$ on $X$, one~has $G=0$ if and only if $f_\ast(G(m))=0$ for $m$ large, for example by virtue of \EGA III 2.2.1 (iii).

\begin{remark} \label{XII.5.6}
One can interpret the criteria \Ref{XII.5.2} and \Ref{XII.5.5} by saying that the \og simultaneous finiteness condition\fg \Ref{XII.5.2} (ii) is expressed by properties of local regularity (in terms of depth thanks to \Ref{VIII}~\Ref{VIII.2.1}) of $F$ at the points of $U$ neighbouring $Y=X- U$, whereas the \og asymptotic vanishing condition\fg \Ref{XII.5.5} (i) is of a distinctly stronger nature, and is expressed by conditions of local regularity of $g_\ast(F)$ at the points of~$Y$ itself. It would be interesting, in order to generalize the Lefschetz-type theorems for projective morphisms to quasi-projective morphisms, to find local criteria on $X$ necessary and sufficient for the $\Ss$-Modules $\bigoplus_{m\geq 0} \R^ih_\ast(F(m))$ for $i\leq n$ to be of finite type. When $S$ is the spectrum of a field (and no doubt more generally, if it is the spectrum of an artinian ring) and $Y=X- U$ is finite, one can show that it is necessary and sufficient that the following conditions be satisfied:
\begin{enumerate}
\item[\textup{1\noo)}] $\prof F_x >n$ for every closed point $x$ of $U$ (compare \Ref{XII.1.4}).
\item[\textup{2\noo)}] $\R^ig_\ast(F)$ is coherent for $i\leq n$, or what amounts to the same, there exists an open neighbourhood $V$ of $Y$ such that for every closed point $x$ of $U\cap V$, one has $\prof F_x >n+1$.
\end{enumerate}
\end{remark}

\begin{proposition} \label{XII.5.7}
Let $S$ be a locally noetherian prescheme, $g:U\to X$ a morphism of preschemes of finite type over $S$\sfootnote{It in fact suffices that $g$ be quasi-compact and quasi-separated (\textup{\EGA IV1.2.1}), with no condition on $U, X$.}, with structural morphisms $h$ and $f$, $F$ a quasi-coherent Module on $U$, $n$ an integer. The following conditions are equivalent:
\begin{enumeratei}
\item
For every base change $S'\to S$, with $S'$ noetherian, the module $\R^ng'_\ast(F')$ on $X'$ is coherent.
\item
For every base change as above, and every coherent Ideal $\Jj$ on~$S'$, denoting by $\Ii$ the Ideal $\Jj\Oo_{X'}$ on $X'$, the graded Module
\[
\tbigoplus_{m\geq 0} \R^ng'_\ast(\Ii^mF')
\]
over $\bigoplus_{m\geq 0} \Ii^m$ is of finite type.
\item
For every base change $S'\to S$, and $\Jj$ as above, the graded Module
\[
\tbigoplus_{m\geq 0} \R^ng'_\ast(\Ii^mF'/\Ii^{m+1}F')
\]
over $\gr_I(\Oo_{X'})=\bigoplus_{m\geq 0} \Ii^m/\Ii^{m+1}$ is of finite type.
\end{enumeratei}
\end{proposition}

Obviously (ii) \ALORS (i) and (iii) \ALORS (i), as one sees by setting $\Ii=0$ in conditions (ii) and (iii). The converse implications are obtained by applying (i) to the composite base change $S^{''}\to S'\to S$, where $S^{''}$ is equal to $\Spec \bigoplus_{m\geq 0} \Jj^m$ \resp $\Spec \bigoplus_{m\geq 0} \Jj^m/\Jj^{m+1}$.

The interest of this proposition is that the conditions of the form (ii) are those which intervene in the \og algebraic-formal comparison theorems\fg, whereas the conditions of the form (iii) intervene in the \og existence theorems\fg which complete them, \cf expos\'e \Ref{IX}. A first interesting case is that where $f:X\to S$ is the identity, and where it is thus a matter of conditions on a morphism $h:U\to S$ locally of finite type and a Module $F$ quasi-coherent on $U$ flat with respect to $S$. To obtain sufficient conditions, we shall assume that $U$ embeds via $g:U\to X$, as an open subprescheme of an $X$ \emph{proper} over $S$. Applying \Ref{XII.5.1}, one thus sees:

\begin{corollary} \label{XII.5.8}
Let $f:X\to S$ be a proper morphism, with $S$ locally noetherian, $U$ an open set of $X$, $g:U\to X$ the canonical immersion, $h=fg:U\to S$, $F$ a quasi-coherent Module on $X$, flat with respect to $S$. Suppose that for every base change $S'\to S$, with $S'$ locally noetherian, one has $\R^ig'_\ast(F')$ coherent on $X'$ for $i\leq n$. Then one~has the following:
\begin{enumeratei}
\item
For every base change $S'\to S$, with $S'$ locally noetherian, $\R^ih'_\ast(F')$ is coherent on $S'$ for $i\leq n$.
\item
For every $S'\to S$ as above, and every coherent Ideal $\Jj$ on $S'$, the graded Modules
\[
\tbigoplus_{m\geq 0} \R^ih'_\ast(\Jj^mF')
\]
over $\bigoplus_{m\geq 0}\Jj^m$ are of finite type for $i\leq n$.
\item
For every $S'\to S$ and $\Jj$ as above, the graded Modules
\[
\tbigoplus_{m\geq 0} \R^ih'_\ast(\Jj^mF'/\Jj^{m+1}F')
\]
over $\gr_\Jj(\Oo_{S'})=\bigoplus_{m\geq 0}\Jj^m/\Jj^{m+1}$ are of finite type for $i\leq n$.

Moreover, under the conditions of (ii), and by virtue of the comparison theorem \Ref{IX.1.1}, denoting by $\widehat{S'}$ the formal completion of $S'$ with respect to $\Jj$, and by $\widehat{U'}$ that of $U'$ with respect to $\Jj\Oo_{U'}$, the canonical homomorphisms
\[
\widehat{\R^ih'_\ast(F')} \to \R^i\widehat{h'}_\ast(\widehat{F'}) \to \varprojlim_k\R^ih'_\ast(F'_k)
\]
are isomorphisms for $i\leq n-1$.
\end{enumeratei}
\end{corollary}

\begin{remark} \label{XII.5.9}
Suppose that one is moreover under the conditions of \Ref{XII.5.8} with $F$ \emph{coherent}, and consider a base change $S'\to S$ as in \Ref{XII.5.9} (i). Suppose moreover that $S'$ is locally immersible in a regular scheme, or more generally, that there exists a morphism $S^{''}\to S$ faithfully flat and quasi-compact, such that $S^{''}$ is locally immersible in a regular scheme; this condition is satisfied in particular if $S'$ is local. Then the conclusion of \Ref{XII.5.8} (i) and (ii) remains valid when one replaces $F'$ therein by a Module $G'$ on $U'$, such that every point of $U'$ has an open neighbourhood on which $G'$ is isomorphic to a Module of the form ${F'}^n$. Indeed, one is reduced to the case $S'$ itself is locally immersible in a regular scheme, so that the same holds for $S^{''}=\Spec\big(\bigoplus_{m\geq 0}\Jj^m\big)$ and for $X^{''}=X'\times_{S'} S^{''}= X\times_{S} S^{''}$, which are of finite type over it. One then applies the finiteness criterion \Ref{VIII}~\Ref{VIII.2.3} for the direct images for $i\leq n$ of $G^{''}$ under the immersion $U^{''}\to X^{''}$, noting that they are satisfied by hypotheses for $F^{''}$, hence also for $G^{''}$ since they are expressed in terms of depth and $G^{''}$ is locally isomorphic to a ${F^{''}}^n$. The same argument shows that if $\Gg'$ is a coherent Module on $\widehat{U'}$ (completion of $U'$ with respect to the Ideal $\Jj\Oo_{U'}$), such that $\Gg'_0=\Gg'/\Jj\Gg'$ is locally of the form $F_0^{'n}$, then the conclusion of (iii) remains valid upon replacing $F'$ therein by $\Gg'$. One thus obtains the following result, using the results of expos\'e \Ref{IX}:
\end{remark}

\begin{corollary} \label{XII.5.10}
Let $f:X\to S$ be a proper morphism, with $S$ locally noetherian, $U$ an open subset of $X$; one assumes $U$ flat with respect to $S$, and that for every base change $S'\to S$, with $S'$ locally noetherian, one has $R^ig'_\ast(\Oo_{U'})$ coherent on $X'$ for $i=0, 1$. Suppose then that $S'$ is of the form $\Spec(A)$, where $A$ is noetherian ring equipped with an ideal $J$ such that $A$ is separated and complete for the $J$-adic topology. Under these conditions:
\begin{enumeratei}
\item
The functor $F\mto\widehat{F}$ from the category of locally free Modules on $U'$ to the category of locally free Modules on $\widehat{U'}$ is fully faithful.
\item
For every locally free Module $\formelF$ on $\widehat{U'}$, there exists a coherent Module $F$ on~$U'$ (not necessarily locally free, alas), and an isomorphism $\widehat{F}\simeq \formelF$.
\end{enumeratei}
\end{corollary}

It remains only to prove (ii), thanks to \Ref{XII.5.9}. Now according to this remark and \Ref{IX.2.1} it follows that $\formelF$ is induced by a coherent Module $\formelG$ on $\widehat{X'}$. By the existence theorem \EGA III 5.1.4, $\formelG$ is of the form $\widehat{F}$, where $F$ is coherent on $X$, whence the conclusion.

\begin{remarks} \label{XII.5.11}
1\textsuperscript{o}) Using \Ref{XII.5.10}, \Ref{XII.4.7} and a suitable hypothesis, saying that certain local rings of the geometric fibres of $X'\to S'$ are \og pure\fg \resp parafactorial, one should be able to obtain statements saying that the functor $Z'\mto \widehat{Z'}$ from the category of \'etale coverings of $X'$ to the category of \'etale coverings of $\widehat{X'}$ (or what amounts to the same, of $X_0'$) is an equivalence of categories, \resp that the functor $L\mto \widehat{L}$ from the category of invertible Modules on $X'$ to the category of invertible Modules on $\widehat{X'}$ is an equivalence. Using recent results of Murre, it is probable that one should be able to deduce existence theorems for Picard schemes for certain \emph{non-proper} algebraic schemes\nde{of course, in the projective case one will refer to Grothendieck's existence theorems of FGA; \cf Grothendieck~A., {\og Technique de descente et th\'eor\`emes d'existence en g\'eom\'etrie alg\'ebrique. VI. Les sch\'emas de Picard: propri\'et\'es g\'en\'erales\fg}, in \emph{S\'eminaire Bourbaki}, vol.~7, Soci\'et\'e math\'ematique de France, Paris, 1995, \Exp 236, p\ptbl 221--243 and {\og Technique de descente et th\'eor\`emes d'existence en g\'eom\'etrie alg\'ebrique. V. Les sch\'emas de Picard: th\'eor\`emes d'existence\fg}, in \emph{S\'eminaire Bourbaki}, vol.~7, Soci\'et\'e math\'ematique de France, Paris, 1995, \Exp 232, 143--161. The nine finiteness conjectures which are found there are proved in expos\'es \Ref{XII} and \Ref{XIII} of Mme Raynaud and of Kleiman of \SGA 6. For an excellent elementary text on the subject, see Kleiman's survey article (Kleiman~S., {\og The Picard scheme\fg}, to appear in Contemp. math.). For an application of these techniques to the global generalized Jacobians of a relative smooth curve, see (Contou-Carr\`ere~C., {\og La jacobienne g\'en\'eralis\'ee d'une courbe relative; construction et propri\'et\'e universelle de factorisation\fg}, \textit{C.~R. Acad. Sci. Paris S\'er.~A-B} \textbf{289}, (1979), \numero3, A203--A206 and {\og Jacobiennes g\'en\'eralis\'ees globales relatives\fg}, in \emph{The Grothendieck Festschrift, Vol. II}, Progr. Math., vol.~87, Birkh\"auser, Boston, 1990, p\ptbl 69--109). See also by the same author, in the purely local setting, the construction and the study of the functor \og local generalized Jacobian\fg ({\og Jacobienne locale, groupe de bivecteurs de Witt universel, et symbole mod\'er\'e\fg}, \textit{C.~R. Acad. Sci. Paris S\'er.~I Math.}  \textbf{318} (1994), \numero 8, p\ptbl 743--746). Moreover, if in the case of a projective and smooth morphism, the connected components of the Picard scheme are proper, this is no longer so in the singular case. There naturally arises the problem of the compactification of Picard schemes: this problem has been studied in detail, notably in (Altman~A.B. \& Kleiman~S., {\og Compactifying the Picard scheme\fg}, \emph{Adv. in Math.} \textbf{35} (1980), \numero 1, p\ptbl 50--112. and {\og Compactifying the Picard scheme. II\fg}, \emph{Amer. J.~Math.} \textbf{101} (1979), \numero 1, p\ptbl 10--41). The case of curves had been studied previously (D'Souza~C., {\og Compactification of generalised Jacobians\fg}, \emph{Proc. Indian Acad. Sci. Sect. A Math. Sci.} \textbf{88} (1979), \numero 5, p\ptbl 419--457). One even knows exactly when the compactified Jacobian of a curve is irreducible (Rego~C.J., {\og The compactified Jacobian\fg} \emph{Ann. Sci. \'Ec. Norm. Sup. (4)} \textbf{13} (1980), \numero 2, p\ptbl 211--223, the latter being the closure of the (ordinary) Jacobian when the curve is geometrically integral and locally planar; for a construction in families of compactified Jacobians, see (Esteves~E., {\og Compactifying the relative Jacobian over families of reduced curves\fg}, \emph{Trans. Amer. Math. Soc.} \textbf{353} (2001), \numero 8, p\ptbl 3045--3095). Since then, the existence results for the Picard scheme in the proper case have progressed since the original edition of \SGA 2; \cf (Murre~J.P., {\og On contravariant functors from the category of pre-schemes over a field into the category of abelian groups (with an application to the Picard functor)\fg}, \emph{Publ. Math. Inst. Hautes \'Etudes Sci.} \textbf{23} (1964), p\ptbl 5-43) and above all (Artin~M., {\og Algebraization of formal moduli. I\fg}, in \emph{Global Analysis (Papers in Honor of K\ptbl Kodaira)}, Univ. Tokyo Press, Tokyo, 1969, p\ptbl 21--71). See also (Raynaud~M., {\og Sp\'ecialisation du foncteur de Picard\fg}, \emph{Publ. Math. Inst. Hautes \'Etudes Sci.} \textbf{38} (1970), p\ptbl 27--76) in the case of a proper scheme over a discrete valuation ring, but not necessarily flat. For a discussion of more recent results, in particular those of Artin, for the Picard functor of proper and flat schemes, in particular in the case cohomologically flat in dimension $0$, see Chapter VIII of (Bosch~S., L\"utkebohmert~W. \& Raynaud~M., \emph{N\'eron models}, Ergebnisse der Mathematik und ihrer Grenzgebiete (3), vol.~21, Springer-Verlag, Berlin, 1990) and the references cited. Much more recently, very fine results have been obtained in the case of relative curves $f:X\to S$ over the spectrum $S$ of a discrete valuation ring with perfect residue field. More precisely, one assumes that $f$ is proper and flat, $X$~regular and $f_*\OX=\Oo_S$. On the other hand, one does not assume $f$ cohomologically flat in dimension~$0$, \ie one does not assume $H^1(X,\Oo)$ torsion-free. The Picard scheme is then not representable, whether by a scheme or an algebraic space, as soon as the torsion in question is non-zero. Let~$J$ be the N\'eron model of the generic fibre of $f$: it is a quotient of the Picard functor $P$. Then, Raynaud has shown that the kernel of the tangent map $H^1(X,\Oo)=\mathrm{Lie}(P)\to \mathrm{Lie}(J)$ coincides with the torsion subgroup of the $H^1$ and that the cokernel has the same length (see theorem 3.1 of (Liu~Q., Lorenzini~D. \& Raynaud~M., {\og N\'eron models, Lie algebras, and reduction of curves of genus one\fg}, \emph{Invent. Math.} \textbf{157} (2004), p\ptbl 455--518). This result allows the aforementioned authors to study the link between the Birch--Swinnerton-Dyer and Artin--Tate conjectures (see th\ptbl 6.6 of {\loccit}). As regards the local Picard scheme, see Boutot's thesis, cited in the editor's note~\eqref{Boutot}~page~\pageref{Boutot}.}. In a general way, the elimination of purity hypotheses in various existence theorems, notably of representability of functors such as the Hilbert functors, or Picard functors etc., with the help of the techniques developed in this seminar, deserves a systematic study.

2\textsuperscript{o}) One may propose to give manageable necessary and sufficient conditions, in terms of depth, for the universal finiteness condition envisaged in \Ref{XII.5.10} to be satisfied. When $S$ is the spectrum of a field, it follows easily from \EGA III 1.4.15 that it is necessary and sufficient that the $\R^ig_\ast(F)$ ($i\leq n$) be coherent, which is well expressed in terms of depth thanks to \Ref{VIII}~\Ref{VIII.2.3}. In the general case, one will note however that it does not suffice to require that the preceding condition be satisfied for all the fibres $U_s\subset X_s$ ($s\in S$), even in the case where $n=0$. Take for example $X=S$, $S$ the spectrum of a discrete valuation ring, $U$ the open set reduced to the generic point, $F=\Oo_U$.

3\textsuperscript{o}) Here however is a \emph{sufficient} condition ensuring that one is under the conditions of the hypothesis of \Ref{XII.5.10}: It suffices that $f$ be flat, and that for every $s\in S$ and every $x\in Y_s=X_s- U_s$, one have
\[
\prof \Oo_{X_s, x}\geq n+2.
\]
Indeed, taking into account lemma \Ref{XII.2.5} (\cf relation (\Ref{eq:XII.16}) after \Ref{XII.2.5}), it follows that one then has $g_\ast(\Oo_U)\simeq\OX$ and $\R^ig_\ast(\Oo_U)=0$ for $0<i\leq n$, and the same relations will obviously be satisfied after every base change $S'\to S$.
\end{remarks}
